\documentclass[11pt,a4paper]{article}
\usepackage[margin=2.2cm]{geometry}
\usepackage{amsmath}

\title{Performance Analysis of LLM Inference Engines on Deucalion CPU Nodes}
\author{Big Data Analysis Project Assignment \#2}
\date{2025/2026}

\begin{document}
\maketitle

\section{Introduction}

This report follows Track A1, Single-Engine Deep Dive. The primary engine is
\texttt{llama.cpp} served through the OpenAI-compatible \texttt{llama-server}
HTTP interface. The experiments also include TurboQuant cache variants,
the vanilla \texttt{llama.cpp} baseline, the same sweep on x86 nodes, and a
vanilla BLAS comparison on ARM. These extra comparisons are used to explain
the same A1 optimisation space against hardware differences rather than to
claim a separate Track A2 study.

The parsed result set contains 1548 prompt-level summary rows and 1548 scaling rows.
The raw sweep directories are \texttt{measurements/1253766-track-a1-server-sweep},
\texttt{measurements/1254633-track-a1-server-sweep-x86}, and
\texttt{measurements/1254637-track-a1-vanilla-blas-sweep}. The organized CSVs,
plots, and this report are under \texttt{measurements/organized\_a1\_results}.

\section{Methodology}

The ARM runs used Deucalion \texttt{normal-arm} nodes with Fujitsu A64FX:
48 online CPUs, four NUMA nodes, one hardware thread per core, and SVE support.
The x86 runs used \texttt{normal-x86} nodes with two AMD EPYC 7742 sockets:
128 online CPUs, eight NUMA nodes, one hardware thread per core, AVX2/FMA, and
256 GB class memory nodes as described in the assignment. The ARM assignment
notes describe 32 GB class memory nodes. These capacity differences matter for
headroom, but decode latency is primarily governed by model streaming and cache
locality rather than capacity alone.

The benchmark client sends streaming HTTP requests and records time to first
token (TTFT), time per output token (TPOT), output-token throughput, goodput,
and sampled peak memory. Each configuration uses three measured trials after
one warmup trial. Goodput uses the harness SLA of TTFT no larger than 2.0 s and
TPOT no larger than 0.2 s. The prompt file contains at least ten prompts per category and
marks the three mandatory assignment prompts. These sweeps used
\texttt{LIMIT\_PER\_CATEGORY=3}, so the measured workload covers nine prompts
per configuration: the three mandatory prompts plus six additional prompts.
That is sufficient for the mandatory-prompt analysis below, but it is a partial
run relative to the strict full-workload requirement of ten measured prompts per
category. Generated text for the mandatory prompts is stored in
\texttt{measurements/organized\_a1\_results/combined/mandatory\_quality.csv}.

\section{Metric Definitions}

TTFT is reported in seconds and is the elapsed time from request submission to
the first streamed output token. It mostly reflects prompt prefill, queueing,
and startup overhead. TPOT is reported in milliseconds per generated token and
is the steady-state decode latency. Throughput is reported as generated output
tokens per second. Peak memory is reported as GiB from sampled VmHWM high-water
resident memory. Goodput is a fraction from 0 to 1 using the benchmark SLA:
TTFT no larger than 2.0 seconds and TPOT no larger than 200 milliseconds.

\section{Results}

\begin{table}[htbp]
\centering
\small
\begin{tabular}{lllllllll}
\hline
Sweep & Engine & Cache & Model & Threads & TTFT s & TPOT ms/token & Output tokens/s & Peak memory GiB \\
\hline
arm\_server & tq & turbo3 & model-1-mandatory & 24 & 0.525 & 145.3 & 5.461 & 4.83 \\
arm\_server & tq & turbo4 & model-1-mandatory & 24 & 0.519 & 140.0 & 5.686 & 4.87 \\
arm\_server & tq & turbo3 & model-3 & 24 & 11.304 & 164.4 & 0.254 & 5.18 \\
arm\_server & tq & turbo4 & model-3 & 24 & 1.585 & 128.8 & 1.158 & 5.19 \\
arm\_server & vanilla & f16 & model-1-mandatory & 24 & 0.555 & 130.4 & 6.130 & 5.25 \\
arm\_server & vanilla & f16 & model-2 & 24 & 0.090 & 19.7 & 40.517 & 0.47 \\
arm\_server & vanilla & f16 & model-3 & 24 & 1.772 & 127.7 & 3.545 & 5.41 \\
blas\_server & vanilla-blis & f16 & model-1-mandatory & 24 & 0.543 & 130.4 & 6.142 & 5.51 \\
blas\_server & vanilla-blis & f16 & model-2 & 24 & 0.089 & 21.1 & 38.120 & 0.51 \\
blas\_server & vanilla-blis & f16 & model-3 & 24 & 1.777 & 127.8 & 3.540 & 5.57 \\
blas\_server & vanilla-openblas-fujitsu & f16 & model-1-mandatory & 24 & 0.555 & 130.6 & 6.125 & 5.47 \\
blas\_server & vanilla-openblas-fujitsu & f16 & model-2 & 24 & 0.088 & 19.6 & 40.679 & 0.50 \\
blas\_server & vanilla-openblas-fujitsu & f16 & model-3 & 24 & 1.777 & 128.5 & 3.526 & 5.54 \\
x86\_server & tq & turbo3 & model-1-mandatory & 32 & 0.191 & 158.9 & 5.924 & 7.99 \\
x86\_server & tq & turbo4 & model-1-mandatory & 32 & 0.160 & 130.6 & 7.186 & 8.03 \\
x86\_server & tq & turbo3 & model-3 & 32 & 0.293 & 93.9 & 2.940 & 7.02 \\
x86\_server & tq & turbo4 & model-3 & 32 & 0.374 & 120.6 & 2.059 & 7.04 \\
x86\_server & vanilla & f16 & model-1-mandatory & 32 & 0.165 & 122.7 & 7.520 & 8.42 \\
x86\_server & vanilla & f16 & model-2 & 32 & 0.040 & 15.5 & 56.248 & 0.50 \\
x86\_server & vanilla & f16 & model-3 & 32 & 0.468 & 129.0 & 4.690 & 7.24 \\
\hline
\end{tabular}
\caption{Baseline mandatory-prompt means with explicit units.}
\label{tab:baseline}
\end{table}

The main baseline pattern is that the x86 node delivers lower TPOT for the
8B mandatory model at the tested baseline, but it also uses a larger resident
memory footprint. On ARM, vanilla \texttt{llama.cpp} and the BLAS-linked vanilla
variants are competitive with or faster than the TurboQuant cache modes for the
tested Q4 model. TurboQuant cache modes are still useful as an optimisation
dimension because they change KV-cache representation, memory footprint, and
decode behaviour, but in these measurements the dominant cost remains full
model-weight streaming during autoregressive decode.

\begin{table}[htbp]
\centering
\small
\begin{tabular}{lllllll}
\hline
Sweep & Engine & Cache & TTFT s & TPOT ms/token & Output tokens/s & Peak memory GiB \\
\hline
arm\_server & tq & turbo3 & 0.525 & 145.3 & 5.461 & 4.83 \\
arm\_server & tq & turbo4 & 0.519 & 140.0 & 5.686 & 4.87 \\
arm\_server & vanilla & f16 & 0.555 & 130.4 & 6.130 & 5.25 \\
x86\_server & tq & turbo3 & 0.191 & 158.9 & 5.924 & 7.99 \\
x86\_server & tq & turbo4 & 0.160 & 130.6 & 7.186 & 8.03 \\
x86\_server & vanilla & f16 & 0.165 & 122.7 & 7.520 & 8.42 \\
\hline
\end{tabular}
\caption{ARM versus x86 baseline comparison. The x86 wrapper used 32 threads; ARM used 24 threads.}
\label{tab:architecture}
\end{table}

Thread scaling shows diminishing returns. On ARM, performance improves when the
thread count rises from one or two CMG-sized groups toward the 36--46 thread
region, but 48 threads is not always best. On x86, 8--16 threads are often
already strong; 128 threads can be worse because the request is spread across
many NUMA domains and synchronization overhead dominates useful matrix-vector
work. Concurrency behaves as expected for CPU serving: increasing simultaneous
requests increases queueing and contention, so per-request TTFT and TPOT rise,
and the goodput SLA degrades. Some high-concurrency BLAS runs returned no
successful prompt rows, which should be treated as overload points rather than
valid latency samples.

Context length mostly affects TTFT and memory because longer prompts increase
prefill work and KV-cache size. Decode length has a smaller effect on TPOT; it
mainly improves measurement stability because fixed startup overhead is spread
over more generated tokens. The small SmolLM2 model is much faster and has a
much smaller memory footprint. The Gemma model produced several zero-token or
weak-output rows, so those model-set results should be discussed as a quality
or compatibility limitation rather than as a clean performance point.

\begin{table}[htbp]
\centering
\small
\begin{tabular}{lllllll}
\hline
Engine & Threads & TTFT s & TPOT ms/token & Output tokens/s & Goodput fraction & Peak memory GiB \\
\hline
vanilla-blis & 24 & 0.543 & 130.4 & 6.142 & 1.000 & 5.51 \\
vanilla-openblas-fujitsu & 24 & 0.555 & 130.6 & 6.125 & 1.000 & 5.47 \\
\hline
\end{tabular}
\caption{BLAS runtime comparison on ARM using vanilla llama.cpp.}
\label{tab:blas}
\end{table}

The BLAS comparison shows OpenBLAS/Fujitsu and BLIS are very close for this
server workload. That is consistent with decode being memory-traffic dominated:
once the matrix-vector kernels are reasonably optimized, changing BLAS does not
remove the need to stream the quantised weights for each output token.

\section{Performance Model}

For autoregressive decoding, a simple lower-bound model is
\[
  \mathrm{TPOT}_{pred} \geq \frac{\mathrm{model\_bytes}}{BW_{mem}}.
\]
The intuition is that each output token requires reading a large fraction of the
model weights. If the run is memory-bandwidth bound, TPOT should scale roughly
with model size and inversely with attainable memory bandwidth. The repository
does not currently contain a \texttt{memory\_bandwidth.csv} STREAM measurement,
so the validation table below uses the best observed effective model-streaming
rate per sweep as an empirical bandwidth ceiling. This is weaker than an
independent STREAM validation and should be replaced by measured TRIAD bandwidth
before final submission.

\begin{table}[htbp]
\centering
\small
\begin{tabular}{llllllll}
\hline
Sweep & Engine & Model & Model GiB & BW GB/s & Pred TPOT ms & Observed TPOT ms & Obs/Pred \\
\hline
arm\_server & tq & model-1-mandatory & 4.58 & 41.793 & 117.7 & 145.3 & 1.23 \\
arm\_server & tq & model-1-mandatory & 4.58 & 41.793 & 117.7 & 140.0 & 1.19 \\
arm\_server & tq & model-3 & 4.97 & 41.793 & 127.7 & 164.4 & 1.29 \\
arm\_server & tq & model-3 & 4.97 & 41.793 & 127.7 & 128.8 & 1.01 \\
arm\_server & vanilla & model-1-mandatory & 4.58 & 41.793 & 117.7 & 130.4 & 1.11 \\
arm\_server & vanilla & model-2 & 0.25 & 41.793 & 6.5 & 19.7 & 3.04 \\
arm\_server & vanilla & model-3 & 4.97 & 41.793 & 127.7 & 127.7 & 1.00 \\
blas\_server & vanilla-blis & model-1-mandatory & 4.58 & 41.751 & 117.9 & 130.4 & 1.11 \\
blas\_server & vanilla-blis & model-2 & 0.25 & 41.751 & 6.5 & 21.1 & 3.25 \\
blas\_server & vanilla-blis & model-3 & 4.97 & 41.751 & 127.8 & 127.8 & 1.00 \\
blas\_server & vanilla-openblas-fujitsu & model-1-mandatory & 4.58 & 41.751 & 117.9 & 130.6 & 1.11 \\
blas\_server & vanilla-openblas-fujitsu & model-2 & 0.25 & 41.751 & 6.5 & 19.6 & 3.02 \\
\hline
\end{tabular}
\caption{Simple TPOT model validation. BW is calibrated from the best observed model-streaming rate in each sweep.}
\label{tab:model}
\end{table}

The model captures the correct first-order trend: the 8B Q4 model is much
slower and much larger than SmolLM2, and the x86 node can sustain lower TPOT at
the tested baseline. It breaks down where engine overhead, NUMA placement,
thread scheduling, prompt processing, KV-cache layout, and failed/zero-token
generations dominate. Ratios above one indicate overhead beyond a pure memory
streaming lower bound. Ratios below or near one identify the rows used to
calibrate the empirical bandwidth ceiling.

\section{Bottleneck Analysis}

The dominant bottleneck is memory traffic during decode, with secondary
bottlenecks from thread coordination and NUMA locality. Evidence:
\begin{itemize}
\item TPOT tracks model size strongly: SmolLM2 has far lower memory use and much higher throughput than the 8B baseline.
\item Decode length changes throughput less than model size or thread placement, meaning steady-state token generation is the limiting phase.
\item High thread counts eventually hurt, especially on x86 at 128 threads, which is typical when memory bandwidth and cross-NUMA traffic saturate.
\item BLAS library changes have small impact compared with architecture, model size, and threading.
\end{itemize}

The highest-impact optimisation would be to reduce bytes moved per generated
token: smaller weight quantisation, better cache locality, NUMA-aware pinning,
and speculative decoding or prompt batching where quality and latency budgets
allow it. KV-cache quantisation helps memory footprint and long-context
pressure, but it cannot fully remove the repeated model-weight traffic in
decode.

\section{Requirement Mapping}

\begin{table}[htbp]
\centering
\small
\begin{tabular}{lll}
\hline
Requirement & Evidence & Status \\
\hline
Track stated & Track A1 stated in scripts, README, and this report & Met \\
At least three models & Llama 3.1 8B Q4, SmolLM2 360M Q4, Gemma Q4 & Met \\
Mandatory baseline model & Meta-Llama-3.1-8B-Instruct-Q4\_K\_M.gguf & Met \\
Metrics & TTFT, TPOT, throughput, goodput, sampled VmHWM memory & Met \\
Three trials plus warmup & \texttt{trials=3}, \texttt{warmup\_trials=1} in system logs & Met \\
Prompt dataset & JSON contains 10+ prompts per category and mandatory prompts & Met in dataset \\
Measured prompt count & Runs used \texttt{LIMIT\_PER\_CATEGORY=3} & Partial \\
Optimisation dimensions & cache type, threads, concurrency, context, decode length, node architecture, BLAS & Met \\
Performance model & TPOT lower bound from model bytes and bandwidth & Partial until STREAM bandwidth is added \\
\hline
\end{tabular}
\caption{Assignment requirement mapping.}
\label{tab:requirements}
\end{table}

\section{Conclusion}

The measurements support a hardware-aware explanation of CPU-only LLM serving:
decode is dominated by repeated movement of model weights, while prefill and
long-context behaviour are reflected more strongly in TTFT and memory. The x86
EPYC node is faster for the tested baseline, but raw thread count is not enough:
oversubscribing across NUMA domains can reduce throughput. On ARM A64FX, using
most but not necessarily all cores is better than assuming maximum thread count
is optimal. BLAS choice matters less than model size, memory traffic, and
thread placement for the measured server workload.

\end{document}
